\uuid{ybdi}
\titre{Raisonnements}
\auteur{bodin ; croizat ; sacre}
\organisation{exo7}
\licence{CC BY-NC-SA 4.0}
\datecreate{2022-03-08}
\niveau{L1}
\module{Algèbre}
\chapitre{Logique, ensemble, raisonnement}
\sousChapitre{Absurde et contraposée}
\difficulte{2}
\theme{Raisonnements}
\exo7qcmid{hilisit.logique.049}

\contenu{
\begin{question}
\qclassification{100.03}
Soient $a$ et $b$ deux nombres réels. On considère la proposition suivante : "si $a+b$ est irrationnel, alors $a$ est irrationnel ou $b$ est irrationnel". Comment puis-je montrer cette affirmation par contraposée ?
\begin{answers}
	\good{Je prends deux rationnels $a$ et $b$ et je montre que $a+b$ est rationnel.}
	\bad{Je prends deux irrationnels $a$ et $b$ et je montre que $a+b$ est irrationnel.}
	\bad{Je prends un irrationnel et j'essaie de l'écrire sous la forme $a+b$ avec $a$ et $b$ irrationnels.}
	\bad{Je prends deux rationnels $a$ et $b$ et je montre que $a+b$ est irrationnel.}
\end{answers}
\begin{explanations}
La contraposée est : si $a$ et $b$ sont rationnels alors $a+b$ est rationnel. On prend donc $a$ et $b$ deux rationnels : $a=\dfrac cd$, $b=\dfrac ef$. Alors $a+b=\dfrac cd+\dfrac ef=\dfrac {cf+de}{df}$ donc $a+b$ est rationnel.
\end{explanations}
\end{question}
}
